% ===== wrap-up =====
\begin{frame}{8주차 정리}
  \begin{enumerate}
    \item \kb{8.1 팀 프로젝트} --- 시드 $\geq 3$ + greedy($\varepsilon=0$)
      \textbf{평가 분리} 프로토콜, 베이스라인, 수식적 정당화(개념+실측
      숫자 삼박자). 채점의 절반 = ``실험 프로토콜의 정직성'' +
      ``결과의 정직성''.
    \item \kb{8.2 Peer-Review} --- 4항목 체크리스트(환경 설명·알고리즘
      선택·수식적 정당화·결과의 정직성) + \textbf{3점 반박 가능한
      질문}(개념 명시 + 답이 새 실험이어야). RL 특유의 ``우연히 좋은''
      패턴 5가지, 특히 패턴 2의 \textbf{비대칭}(1차 vs 2차 추락 확률).
    \item \kb{8.3 총정리} --- 새 개념 0개. 손유도 템플릿 3종
      (거꾸로 대입 / 후보+검증 / 축소 사상 수렴) + 개념 지도의
      \textbf{엣지가 시험}.
  \end{enumerate}
  \vskip0.8em
  \begin{block}{다음 주 --- Week 9. DQN}
    이 주에 체득한 \textbf{Q-함수 · 벨만 최적방정식 · on/off-policy}가
    곧 DQN의 뼈대 --- Q-테이블을 \textbf{신경망}으로 통째로 바꾸고,
    타겟 네트워크와 경험 재현으로 표 기반에서 불가능했던 \textbf{연속
    상태 공간}으로 나아간다. 시드·평가 분리·반박 가능한 질문의 습관이
    Block B 캡스톤까지 그대로 이어진다.
  \end{block}
\end{frame}
